%! Zeros of Polynomials: the Rational Root Theorem — Unit 3, Lesson 3.4 — Guided Notes (Answer Key)
\documentclass[10pt]{article}
\usepackage{algebra2-article}
\usepackage{algebra2-key}   % pulls in -boxes; do NOT also load -boxes
\newcommand{\TallMath}[1]{$\displaystyle #1\rule[-1.4em]{0pt}{3.2em}$}
\newcommand{\mans}[1]{{\color{keyred}\boldsymbol{#1}}}   % math-mode answer (\ans is text-only)

\newcommand{\lgrid}[4]{%
  \draw[step=1, linegray!40, very thin] (#1,#3) grid (#2,#4);
  \draw[->, semithick, charcoal] (#1,0)--(#2,0) node[below right, font=\scriptsize]{$x$};
  \draw[->, semithick, charcoal] (0,#3)--(0,#4) node[left, font=\scriptsize]{$y$};}

\begin{document}

\pageheader{Unit 3, Lesson 3.4}{Guided Notes --- Answer Key}
\namedateperiod

\vspace{0.08in}

\begin{objectivebox}
\textbf{By the end of this lesson, I will be able to\dots}
\begin{itemize}\setlength{\itemsep}{2pt}
  \item use the \emph{Rational Root Theorem} to list every possible rational zero of a polynomial with integer coefficients;
  \item test candidates efficiently with synthetic division, a graph, and the Remainder Theorem;
  \item divide out each zero I find, factor the depressed polynomial, and write the polynomial \emph{completely} factored;
  \item solve a polynomial equation of degree $3$ or higher and \emph{verify} my solutions algebraically and against a graph.
\end{itemize}
\end{objectivebox}

\vspace{0.05in}

\begin{vocabbox}
\small
Fill in each term as we build it together.
\vspace{2pt}

\termblanklong{Rational Root Theorem}
\ansline{If $\frac{p}{q}$ (in lowest terms) is a zero of a polynomial with integer coefficients, then $p$ divides the constant term and $q$ divides the leading coefficient.}
\termblanklong{Candidate (possible rational zero)}
\ansline{One of the values $\pm\frac{p}{q}$ on the list --- a number worth testing, not a guaranteed zero.}
\termblanklong{Depressed polynomial}
\ansline{The quotient left after dividing out a factor $(x-r)$ --- one degree lower, with a shorter candidate list of its own.}
\termblanklong{Solution / root of an equation}
\ansline{A value of $x$ that makes $f(x)=0$ true --- the same thing as a zero of $f$ and an $x$-intercept of its graph.}
\termblanklong{Verifying a solution}
\ansline{Substituting it back to get $0$, and confirming it against the graph's $x$-intercepts.}
\end{vocabbox}

\begin{hookbox}
Yesterday you could finish any polynomial \emph{once somebody handed you a zero to test}. Here is
the polynomial from the Warm-Up, $h(x) = 2x^3 - 3x^2 - 11x + 6$, with its graph. (The $y$-axis is
compressed: each vertical gridline is $4$ units.)
\begin{center}
\begin{tikzpicture}[scale=0.5]
  \lgrid{-3}{4}{-4}{4}
  \begin{scope}
    \clip (-3,-4) rectangle (4,4);
    \draw[very thick, forest, domain=-2.85:3.6, samples=150, smooth]
      plot (\x,{(2*(\x)*(\x)*(\x) - 3*(\x)*(\x) - 11*(\x) + 6)/4});
  \end{scope}
  \foreach \x in {-2,-1,1,2,3}
    \draw[charcoal] (\x,-0.15)--(\x,0.15) node[below=1pt, font=\tiny]{$\x$};
  \fill[charcoal] (-2,0)  circle (2.5pt);
  \fill[charcoal] (0.5,0) circle (2.5pt);
  \fill[charcoal] (3,0)   circle (2.5pt);
\end{tikzpicture}
\end{center}
\begin{itemize}\setlength{\itemsep}{1pt}
  \item Read the three $x$-intercepts off the graph: \ans{$x = -2,\ \tfrac12,\ 3$}
  \item Your group's candidate list yesterday was the divisors of $6$. Which zero did it miss?
        \ans{$\tfrac12$}
  \item That zero is a \emph{fraction}. So a candidate list made only of whole numbers is
        \ans{incomplete}.
\end{itemize}
\textbf{Today's question: which numbers are worth testing?}
\end{hookbox}

\begin{notesbox}{1. Why Guessing Cannot Work}
Suppose nobody tells you anything about $f(x) = x^3 - 4x^2 + x + 6$ except that it has integer
coefficients, and you want its zeros.
\begin{itemize}[leftmargin=1.5em, itemsep=3pt]
  \item How many numbers are there to try? \ans{infinitely many}
  \item So testing candidates one at a time is only useful if the list of candidates is
        \ans{finite (short)}.
\end{itemize}
\textbf{The Rational Root Theorem does exactly that:} it turns two coefficients --- the constant term
and the leading coefficient --- into a \emph{short, finite} list that is guaranteed to contain every
rational zero the polynomial has.
\end{notesbox}

\begin{notesbox}{2. The Rational Root Theorem}
\begin{tcolorbox}[colback=white, colframe=navy, boxrule=0.7pt, arc=1.5mm,
    left=2mm, right=2mm, top=1.5mm, bottom=1.5mm]
\textbf{Rational Root (Zero) Theorem.} Let $f(x) = a_nx^n + \cdots + a_1x + a_0$ have
\emph{integer} coefficients. If $\dfrac{p}{q}$ is a rational zero of $f$ written in lowest terms,
then
\begin{itemize}[leftmargin=1.5em, nosep]
  \item $p$ (the numerator) is a divisor of the \ans{constant term} $a_0$, and
  \item $q$ (the denominator) is a divisor of the \ans{leading coefficient} $a_n$.
\end{itemize}
\end{tcolorbox}

\vspace{4pt}
\textbf{Why it is true --- look at the Warm-Up.} You factored $h(x) = 2x^3-3x^2-11x+6$ into
$(x-3)(2x-1)(x+2)$. Multiply the pieces back:
\begin{itemize}[leftmargin=1.5em, itemsep=3pt]
  \item The \emph{leading} coefficients of the three factors are $1$, $2$, $1$. Their product is
        \ans{$2$}, which is the leading coefficient of $h$. So the $2$ in the denominator of the
        zero $\tfrac12$ had to divide \ans{$2$}.
  \item The \emph{constant} terms of the three factors are $-3$, $-1$, $2$. Their product is
        \ans{$6$}, which is the constant term of $h$. So every numerator had to divide
        \ans{$6$}.
\end{itemize}
That is the whole theorem: a zero $\frac{p}{q}$ comes from a factor $(qx - p)$, and those $q$'s and
$p$'s have nowhere to hide --- they multiply together to make the two coefficients on the ends.

\vspace{5pt}
\textbf{The warnings.} \; (1) It only finds \ans{rational} zeros --- irrational and imaginary zeros
never appear on the list. \; (2) The coefficients must be \ans{integers}. \; (3) The list is a list of
\emph{possibilities}; a polynomial might have \ans{none} of them as an actual zero.
\end{notesbox}

\begin{notesbox}{3. Building the List}
Write $\pm\dfrac{\text{divisors of the constant term}}{\text{divisors of the leading coefficient}}$,
reduce, and cross out repeats.

\vspace{5pt}
\textbf{(a)} $f(x) = x^3 - 4x^2 + x + 6$

\vspace{2pt}
$p$ (divisors of \ans{$6$}): \ans{$\pm 1, \pm 2, \pm 3, \pm 6$} \qquad
$q$ (divisors of \ans{$1$}): \ans{$\pm 1$}

\vspace{2pt}
Candidates $\dfrac{p}{q}$: \ans{$\pm 1,\ \pm 2,\ \pm 3,\ \pm 6$ \; (8 of them)}

\vspace{2pt}
\textbf{Notice:} when the leading coefficient is $1$, every $q$ is $\pm 1$, so the candidates are
just the \ans{divisors of the constant term}. That covers most problems you will see.

\vspace{6pt}
\textbf{(b)} $g(x) = 4x^3 - 4x^2 - x + 1$

\vspace{2pt}
$p$ (divisors of \ans{$1$}): \ans{$\pm 1$} \qquad
$q$ (divisors of \ans{$4$}): \ans{$\pm 1, \pm 2, \pm 4$}

\vspace{2pt}
Candidates $\dfrac{p}{q}$: \ans{$\pm 1,\ \pm\tfrac12,\ \pm\tfrac14$} \qquad How many? \ans{$6$}

\vspace{2pt}
\textbf{The error to avoid:} the constant term goes on \ans{top} and the leading coefficient
goes on \ans{the bottom}. Flipping them is the mistake of the day.
\end{notesbox}

\begin{notesbox}{4. The Workflow: List $\rightarrow$ Test $\rightarrow$ Divide $\rightarrow$ Finish}
\textbf{Model.} Find all the zeros of $f(x) = x^3 - 4x^2 + x + 6$.

\vspace{4pt}
\textbf{Step 1 --- List.} From part (a): \ans{$\pm 1,\ \pm 2,\ \pm 3,\ \pm 6$}

\vspace{4pt}
\textbf{Step 2 --- Test} (Remainder Theorem, or a synthetic division). \;
$f(1) = $ \ans{$4$} \quad $f(-1) = $ \ans{$0$}

\vspace{2pt}
So the first zero is $r = $ \ans{$-1$}, and $($\ans{$x+1$}$)$ is a factor.

\vspace{4pt}
\textbf{Step 3 --- Divide it out.}
\begin{center}
\renewcommand{\arraystretch}{1.6}
\begin{tabular}{r|C{1.4cm}C{1.4cm}C{1.4cm}C{1.4cm}}
\ans{$-1$} & $1$ & $-4$ & $1$ & $6$ \\
    &     & \ans{$-1$} & \ans{$5$} & \ans{$-6$} \\ \cline{2-5}
    & \ans{$1$} & \ans{$-5$} & \ans{$6$} & \ans{$0$} \\
\end{tabular}
\end{center}
Depressed polynomial: \ans{$x^2 - 5x + 6$} \quad Remainder: \ans{$0$}

\vspace{4pt}
\textbf{Step 4 --- Finish.} The depressed polynomial is a quadratic, so no more guessing is needed
--- factor it (or use the quadratic formula).

\vspace{2pt}
$f(x) = $ \ans{$(x+1)(x-2)(x-3)$} \qquad Zeros: \ans{$-1,\ 2,\ 3$}

\vspace{4pt}
\textbf{Step 5 --- Verify} (this is part of the standard, not an extra). Check one zero by
substituting: $f(3) = $ \ans{$27-36+3+6 = 0$}, and check all three against the degree --- a cubic has at most
\ans{$3$} zeros, and you found \ans{$3$}.
\end{notesbox}

\begin{notesbox}{5. Three Ways to Make the Search Shorter}
Twelve candidates is a lot of synthetic divisions. You almost never need them all.

\vspace{4pt}
\textbf{(i) Use the graph.} Here is $f(x) = x^3 - 4x^2 + x + 6$ from Section 4. (Each vertical
gridline is $2$ units.)
\begin{center}
\begin{tikzpicture}[scale=0.5]
  \lgrid{-2}{4}{-3}{4}
  \begin{scope}
    \clip (-2,-3) rectangle (4,4);
    \draw[very thick, forest, domain=-1.7:3.9, samples=150, smooth]
      plot (\x,{((\x)*(\x)*(\x) - 4*(\x)*(\x) + (\x) + 6)/2});
  \end{scope}
  \foreach \x in {-1,1,2,3}
    \draw[charcoal] (\x,-0.15)--(\x,0.15) node[below=1pt, font=\tiny]{$\x$};
\end{tikzpicture}
\end{center}
Which three candidates would you test \emph{first}, and why? \ansline{$-1$, $2$, and $3$ --- the graph crosses the $x$-axis at those three values, and all three are on the candidate list.}

\vspace{4pt}
\textbf{(ii) Shrink the problem.} Once one zero is found, work with the \ans{depressed} polynomial ---
its degree is one lower, and its own candidate list is \ans{shorter}. When it becomes a quadratic,
stop hunting and \ans{factor it or use the quadratic formula}.

\vspace{4pt}
\textbf{(iii) Test cheaply.} The Remainder Theorem evaluates a candidate in a few seconds. Rule out
the losers by \ans{substituting $r$} and save synthetic division for the winners.
\end{notesbox}

\begin{notesbox}{6. Solving a Polynomial Equation}
``Find the zeros of $f$'' and ``solve $f(x) = 0$'' are the \emph{same task} in different words. The
solutions of the equation are the \ans{zeros} of the function, which are the
\ans{$x$-intercepts} of its graph.

\vspace{5pt}
\textbf{Together:} Solve $2x^3 + 3x^2 - 8x + 3 = 0$.

\vspace{2pt}
Candidates: $p$ divides \ans{$3$}, $q$ divides \ans{$2$}, so
$\dfrac{p}{q} = $ \ans{$\pm 1,\ \pm 3,\ \pm\tfrac12,\ \pm\tfrac32$}

\vspace{3pt}
Test $x = 1$: \ans{$2+3-8+3 = 0$} \; so $(x-1)$ is a factor.
\begin{center}
\renewcommand{\arraystretch}{1.6}
\begin{tabular}{r|C{1.4cm}C{1.4cm}C{1.4cm}C{1.4cm}}
\ans{$1$} & $2$ & $3$ & $-8$ & $3$ \\
    &     & \ans{$2$} & \ans{$5$} & \ans{$-3$} \\ \cline{2-5}
    & \ans{$2$} & \ans{$5$} & \ans{$-3$} & \ans{$0$} \\
\end{tabular}
\end{center}
Depressed polynomial: \ans{$2x^2 + 5x - 3$} \; which factors as \ans{$(2x-1)(x+3)$}

\vspace{2pt}
Completely factored: $2x^3+3x^2-8x+3 = $ \ans{$(x-1)(2x-1)(x+3)$}

\vspace{2pt}
Solution set: \ans{$\{\,-3,\ \tfrac12,\ 1\,\}$} \; --- and notice one solution is a fraction, which the divisors of
$3$ alone would never have produced.
\end{notesbox}

\begin{practicebox}
\textbf{A. Run the whole workflow.} Solve $2x^3 - 9x^2 + 7x + 6 = 0$.
\begin{enumerate}[leftmargin=1.7em, itemsep=6pt]
  \item Candidates: \ans{$\pm 1,\ \pm 2,\ \pm 3,\ \pm 6,\ \pm\tfrac12,\ \pm\tfrac32$}
  \item A zero that works: $r = $ \ans{$2$} \quad (show your test: \ans{$16-36+14+6=0$})
  \item Depressed polynomial after dividing: \ans{$2x^2 - 5x - 3$}
  \item Completely factored: \ans{$(x-2)(2x+1)(x-3)$}
  \item Solutions: \ans{$x = -\tfrac12,\ 2,\ 3$}
\end{enumerate}

\vspace{5pt}
\textbf{B. The word ``rational'' is doing real work.} Let $k(x) = x^3 - x^2 - 2x + 2$.
\begin{enumerate}[leftmargin=1.7em, itemsep=6pt, start=6]
  \item Candidates: \ans{$\pm 1,\ \pm 2$} \quad Which one is a zero? \ans{$1$}
  \item Divide it out. Depressed polynomial: \ans{$x^2 - 2$}
  \item Solve that quadratic: $x = $ \ans{$\pm\sqrt{2}$} \quad Are those on your candidate list?
        \ans{No}
  \item How many real zeros does $k$ have in all? \ans{$3$} \; Explain why the Rational Root
        Theorem could never have listed two of them. \ansline{$\pm\sqrt2$ are irrational --- they cannot be written as a fraction $\frac{p}{q}$ of integers, and the theorem only ever produces rational candidates.}
  \item \textbf{Looking ahead.} Suppose a different cubic's depressed polynomial had come out
        $x^2 + 4$. How many \emph{real} zeros would that quadratic contribute? \ans{$0$} \;
        What kind of numbers would its solutions be? \ans{imaginary ($\pm 2i$)} \; That is Lesson 3.5.
\end{enumerate}
\end{practicebox}

\begin{teachernote}
\textbf{Pacing.} Sections 1--3 are the theorem and the bookkeeping; budget about fifteen minutes and
do not rush Section 3, because \emph{building the list correctly} is what most students get wrong on
the test. Sections 4 and 6 are the same workflow twice --- once for ``find the zeros,'' once for
``solve the equation'' --- and the point of doing both is that students recognize them as one task.
If time runs short, cut Section 5(iii), never Section 5(i).

\vspace{3pt}
\textbf{The three errors to pre-empt.} (1) \emph{The fraction upside down} --- constant over leading
coefficient, always. A student who lists $\pm1,\pm2,\pm4$ for $g(x)=4x^3-4x^2-x+1$ has flipped it;
ask them to check whether $4$ could possibly be a zero (it is not remotely close). (2) \emph{Losing
the $\pm$} --- half the list disappears, and in Section 4 the only zero found by testing is negative.
(3) \emph{Duplicates and unreduced fractions} --- $\tfrac22$ and $\tfrac63$ are already on the list;
this is why the Warm-Up made them reduce by hand.

\vspace{3pt}
\textbf{The moment that matters} is Section 2's ``why.'' Students accept the theorem as a rule, but
the one-minute argument makes it inevitable: a zero $\frac{p}{q}$ means a factor $(qx-p)$, all the
$q$'s multiply to the leading coefficient and all the $p$'s multiply to the constant term, so
neither has anywhere to hide. Run it on the Warm-Up's $(x-3)(2x-1)(x+2)$ live, on the board.

\vspace{3pt}
\textbf{Practice B is the bridge to Lesson 3.5} and should not be cut. Item 8's $\pm\sqrt2$ makes
``rational'' a real restriction rather than a word in the theorem's name, and item 10 puts $x^2+4$ on
the table one day before the Fundamental Theorem of Algebra explains what to do with it. Leave item
10 unresolved --- students only need to notice that the depressed quadratic can run out of real
answers.
\end{teachernote}

\end{document}
